% GENEROVÁNO skriptem doc/tools/md2tex.py ze souboru kapitoly/06_vysledky_a_srovnani.md. Needitovat, upravuje se Markdown.

% Kapitola 6, verze 1 (13. 11. 2026): tabulky generuje doc/tools/tabulky_06.py z CSV (bloky mezi  TAB  a  /TAB  needitovat, po novém měření spustit skript znovu). Text je první návrh podle osnovy v doc/osnova_kap4-6.md, čísla z DP/results/vysledky.md (V#) a z CSV běhu na desce 25. 10. 2026 22:43.

\chapter{Výsledky a~srovnání}
\label{kap:vysledky}

Kapitola odpovídá na 4. bod zadání: v~čem se liší jádra napsaná v~HLS a~ve VHDL a~o~kolik je hardwarové zpracování rychlejší než software. Postupuji od správnosti přes prostředky a~maximální frekvenci jednotlivých jader k~propustnosti na desce a~ke srovnání se softwarem. Metodiku měření popisuje kapitola~\ref{kap:metodika}.

\section{Ověření funkce}
\label{sec:vfunkce}

Jádra jsem ověřil na třech úrovních popsaných v~oddílu~\ref{sec:overeni}. Nejdřív jsem porovnal referenční model s~knihovnou OpenCV, poté jádra HLS i~VHDL v~simulaci s~referenčním modelem a~nakonec obě verze na desce. Výsledky shrnuje tabulka~\ref{tab:overeni}.

% TAB tab:overeni

\begin{table}[htbp]
    \centering
    \caption{Ověření funkce: model proti OpenCV, simulace jader proti modelu a~měření na desce (počet úrovní -O0, -O2, -O3 s~bitovou shodou)}
    \label{tab:overeni}
    \begin{tabular}{p{3.4cm}p{3.6cm}p{3.2cm}p{1.7cm}p{1.7cm}}
        \toprule
        Algoritmus & Model proti OpenCV & Simulace HLS i~VHDL & Deska HLS & Deska VHDL \\
        \midrule
        převod barev RGB→YCbCr & odchylka nejvýš 1 & PASS & 3/3 & 3/3 \\
        RGB→HSV & bitově shodný (2^24 barev) & PASS, úplný test 2^24 & 3/3 & 3/3 \\
        HSV→RGB & odchylka nejvýš 2 & PASS, úplný test 2^24 & 3/3 & 3/3 \\
        prahování (Otsu) & bitově shodný & PASS & 3/3 & 3/3 \\
        medián 3×3 & bitově shodný & PASS & 3/3 & 3/3 \\
        Sobel X 3×3 & bitově shodný & PASS & 3/3 & 3/3 \\
        histogram & neporovnáno & PASS & 3/3 & 3/3 \\
        \bottomrule
    \end{tabular}
\end{table}

% /TAB

Na desce jsem změřil všech 7 algoritmů v~obou verzích a~2 průchody řetězcem bez jádra, tedy 16 měření. Každé měření jsem zopakoval s~měřicí aplikací přeloženou na úrovních optimalizace -O0, -O2 a~-O3. Ve všech 48 případech byl výstup jádra bitově shodný s~referenčním modelem přeloženým pro procesor (0 chybných bajtů, u~histogramu 0 chybných košů).

Bitová shoda s~referenčním modelem neznamená bitovou shodu s~OpenCV. U~převodu RGB na šedou a~na HSV, prahování, mediánu a~konvoluce se model s~OpenCV shoduje bit po bitu. U~převodu RGB na YCbCr a~zpět a~u~převodu HSV na RGB se liší nejvýš o~1 až 2 úrovně z~256. OpenCV zde počítá v~plovoucí řádové čárce nebo s~jinak zaokrouhlenými koeficienty, zatímco hardware používá pevnou řádovou čárku s~omezenou šířkou slova. Hardwarové jádro je tedy aproximací, u~které jsem zvolil kompromis mezi přesností, rychlostí a~velikostí obvodu. Například zpětný převod YCbCr na RGB se od OpenCV liší nejvýš o~1 úroveň u~25,5\,\% složek. Bitovou shodu by dal offset s~15 bity za řádovou čárkou, který by ale vyžadoval 26bitový registr místo 10bitového a~jiné jádro než pro dopředný převod. Odchylka o~1 úroveň z~256 odpovídá kvantizaci na 8 bitů a~okem ji v~obraze nerozeznám.

% TODO(Karel): poslední věta je tvoje formulace z 23. 10. („pro lidské oko prakticky nepostřehnutelné“). Chceš ji nechat jako vlastní pozorování, nebo ji podložit (např. rozdílovým obrázkem nebo odkazem na práh rozlišitelnosti)?

Zpětný převod YCbCr na RGB jsem ověřil v~simulaci (xsim a~kosimulace, předvolby 6 až 8), měřicí aplikace ho na desce zatím netestuje.

\section{Prostředky a~maximální frekvence jader}
\label{sec:vooc}

Prostředky a~nejkratší periodu jsem měřil syntézou mimo kontext (oddíl~\ref{sec:ooc}). Pro jádra z~HLS i~VHDL jsem použil stejný skript, obvod xck26-sfvc784-2LV-c, cíl 5\,ns a~Vivado 2025.2. Perioda je větší z~hodnot 5\,ns − WNS a~5\,ns − WPWS, takže zahrnuje i~minimální šířku pulzu hodin bloků BRAM (3,077\,ns) a~DSP (3,262\,ns). Tabulka~\ref{tab:ooc} porovnává celá jádra včetně řízení přes AXI4-Lite, tedy ve stejné podobě, v~jaké jsou v~testovacím řetězci.

% TAB tab:ooc

\begin{table}[htbp]
    \centering
    \caption{Prostředky a~nejkratší perioda jader HLS a~VHDL (syntéza mimo kontext, xck26-sfvc784-2LV-c, cíl 5\,ns, Vivado 2025.2, jádro včetně řízení AXI4-Lite)}
    \label{tab:ooc}
    \begin{tabular}{p{2.6cm}p{1.2cm}p{1.1cm}p{1.1cm}p{1.0cm}p{1.2cm}p{1.4cm}p{1.3cm}p{2.8cm}}
        \toprule
        Jádro & Impl. & LUT & FF & DSP & BRAM36 & Perioda [ns] & Fmax [MHz] & Omezení \\
        \midrule
        převod barev 3→3 & HLS & 407 & 361 & 9 & 0 & 3,88 & 258 & logika \\
         & VHDL & 210 & 389 & 9 & 0 & 3,26 & 307 & šířka pulzu DSP \\
        RGB→HSV & HLS & 249 & 214 & 2 & 1 & 4,44 & 225 & logika \\
         & VHDL & 275 & 255 & 2 & 0 & 3,26 & 307 & šířka pulzu DSP \\
        HSV→RGB & HLS & 556 & 296 & 1 & 0 & 4,05 & 247 & logika \\
         & VHDL & 603 & 422 & 1 & 0 & 3,54 & 283 & logika \\
        prahování & HLS & 122 & 118 & 0 & 0 & 2,77 & 360 & šířka pulzu hodin \\
         & VHDL & 81 & 141 & 0 & 0 & 2,77 & 360 & šířka pulzu hodin \\
        medián 3×3 & HLS & 701 & 472 & 0 & 2 & 3,99 & 251 & logika \\
         & VHDL & 643 & 532 & 0 & 2 & 4,31 & 232 & logika \\
        konvoluce 3×3 & HLS & 640 & 621 & 10 & 2 & 4,05 & 247 & logika \\
         & VHDL & 508 & 558 & 14 & 2 & 3,57 & 280 & logika \\
        konvoluce 5×5 & HLS & 950 & 861 & 26 & 5 & 4,49 & 223 & logika \\
         & VHDL & 766 & 951 & 38 & 4 & 3,76 & 266 & logika \\
        konvoluce 7×7 & HLS & 1 341 & 1 247 & 50 & 7 & 4,35 & 230 & logika \\
         & VHDL & 1 232 & 1 489 & 72 & 6 & 4,00 & 250 & logika \\
        histogram & HLS & 300 & 422 & 0 & 1 & 3,08 & 325 & šířka pulzu BRAM \\
         & VHDL & 183 & 288 & 0 & 1 & 3,08 & 325 & šířka pulzu BRAM \\
        \bottomrule
    \end{tabular}
\end{table}

% /TAB

% Data: HLS řádková jádra (převod barev, RGB→HSV, HSV→RGB) jsou bez auto-rewind, tedy stejné RTL jako v bitstreamu (fronta final_hls_rows 16. 11. 2026, vysledky.md sekce 21).

Verze ve VHDL má kratší nebo stejnou periodu u~8 jader z~9. Největší rozdíl je u~převodu RGB na HSV (3,26\,ns proti 4,44\,ns), konvoluce 5×5 (3,76\,ns proti 4,49\,ns) a~převodu barev (3,26\,ns proti 3,88\,ns). U~převodu RGB na HSV a~převodu barev dosáhlo jádro ve VHDL minimální periody bloku DSP, takže ho dál omezuje obvod, ne logika. Prahování a~histogram mají v~obou verzích stejnou periodu danou minimální šířkou pulzu hodin. Výjimkou je medián, kde má VHDL periodu 4,31\,ns a~HLS 3,99\,ns. Kritická cesta VHDL ale neleží v~síti mediánu, ale v~řízení generátoru okna (příznaky pozice v~řádku a~sloupci), a~70\,\% jejího zpoždění tvoří propojení. Perioda této obálky se mezi běhy měnila od 3,86 do 4,31\,ns, přestože se generátor okna neměnil.

% TODO(Karel): zvážit optimalizaci window_gen (příznaky pozice předpočítat o takt dřív, poznámka z 15. 10.) – zlepšila by medián i konvoluci; jinak text nechat takto.  Příčina je ve způsobu, jakým obě verze vkládají registry. HLS přidá stupňů pipeline jen tolik, aby splnil zadaných 5 ns (RGB→HSV 7 stupňů), ve VHDL jsem dal registr za každou operaci (8 stupňů). Méně registrů znamená méně FF, ale delší logiku mezi nimi.

V~počtu \acs{LUT} je verze ve VHDL menší u~7 jader ze 9. Výjimkou jsou oba převody mezi RGB a~HSV, kde je HLS menší o~9\,\% a~8\,\% LUT a~o~16\,\% a~30\,\% FF. U~převodu barev zabírá VHDL 210 LUT proti 407, u~histogramu 183 proti 300. Konvoluce ve VHDL zabírá méně LUT, ale víc bloků DSP (u~jádra 7×7 je to 72 proti 50), protože jsem každé násobení umístil do vlastního bloku DSP včetně sčítání v~jeho kaskádě. HLS část součinů s~malými koeficienty počítá v~LUT.

\subsection{Jeden parametr: násobení v~LUT, nebo v~DSP}

Převod HSV na RGB obsahuje pět násobení 8\,×\,8 bitů. Obě verze je bez dalšího nastavení počítají v~LUT. Jak se jádro změní, když násobení přesunu do bloků DSP, ukazuje tabulka~\ref{tab:hsvdsp}.

% TAB tab:hsvdsp

\begin{table}[htbp]
    \centering
    \caption{Převod HSV→RGB: násobení v~LUT nebo v~blocích DSP (syntéza mimo kontext, cíl 5\,ns)}
    \label{tab:hsvdsp}
    \begin{tabular}{p{1.3cm}p{4.0cm}p{1.4cm}p{1.1cm}p{1.1cm}p{1.0cm}p{1.4cm}p{2.8cm}}
        \toprule
        Impl. & Parametr & Stupně & LUT & FF & DSP & Perioda [ns] & Omezení \\
        \midrule
        VHDL & \texttt{G\_USE\_DSP = false} & 13 & 603 & 422 & 1 & 3,54 & logika \\
        VHDL & \texttt{G\_USE\_DSP = true} & 13 & 203 & 235 & 11 & 3,26 & šířka pulzu DSP \\
        HLS & výchozí & 6 & 556 & 296 & 1 & 4,05 & logika \\
        HLS & \texttt{HSV\_MUL\_IMPL\_DSP = 1} & 6 & 224 & 241 & 11 & 4,23 & logika \\
        HLS & \texttt{HSV\_MUL\_IMPL\_DSP = 2} (latence 3) & 10 & 262 & 273 & 6 & 3,26 & šířka pulzu DSP \\
        \bottomrule
    \end{tabular}
\end{table}

% /TAB

Ve VHDL stačil jeden parametr (\texttt{G\_USE\_DSP}). Počet LUT klesl o~66\,\%, FF o~44\,\% a~perioda se zkrátila z~3,54 na 3,26\,ns, což je minimální perioda bloku DSP. Cenou je 10 bloků DSP navíc. Pipeline je ve VHDL navržená ručně, takže syntéza mohla registry přesunout do vnitřních registrů bloku DSP.

V~HLS jsem stejnou změnu udělal pragmatem \texttt{BIND\_OP} s~implementací v~DSP. Plocha klesla podobně (LUT o~60\,\%, FF o~19\,\%), ale perioda se prodloužila z~4,05 na 4,23\,ns. HLS ponechal 6 stupňů pipeline a~blok DSP použil jako kombinační násobičku bez vnitřních registrů. Stejné periody jako ve VHDL jsem v~HLS dosáhl až druhým parametrem, latencí násobení 3 takty. Pipeline se tím prodloužila na 10 stupňů a~HLS spojil část násobení do 6 bloků DSP. Ve VHDL tedy stačí jeden parametr, v~HLS jsou potřeba dva a~vývojář musí vědět, že druhý existuje.

Stejný parametr má i~konvoluce ve VHDL (tabulka~\ref{tab:filtrdsp}). U~jádra 7×7 s~násobením v~LUT zabírá 5 238 LUT a~2 702 FF, s~násobením v~blocích DSP 1 213 LUT, 837 FF a~72 bloků DSP. Perioda se změnila nejvýš o~0,4\,ns, u~jádra 3×3 a~7×7 se s~bloky DSP mírně prodloužila. U~konvoluce tedy přesun do DSP šetří hlavně plochu.

% TAB tab:filtrdsp

\begin{table}[htbp]
    \centering
    \caption{Konvoluce K×K ve VHDL: násobení v~LUT nebo v~blocích DSP (jádro bez AXI4-Lite, syntéza mimo kontext, cíl 5\,ns)}
    \label{tab:filtrdsp}
    \begin{tabular}{p{1.4cm}p{1.6cm}p{1.3cm}p{1.3cm}p{1.1cm}p{1.3cm}p{1.6cm}p{2.8cm}}
        \toprule
        K & Násobení & LUT & FF & DSP & BRAM36 & Perioda [ns] & Omezení \\
        \midrule
        3×3 & LUT & 1 120 & 679 & 1 & 2 & 3,26 & šířka pulzu DSP \\
         & DSP & 350 & 318 & 14 & 2 & 3,55 & logika \\
        5×5 & LUT & 2 789 & 1 510 & 1 & 4 & 4,20 & logika \\
         & DSP & 603 & 539 & 38 & 4 & 3,82 & logika \\
        7×7 & LUT & 5 238 & 2 702 & 1 & 6 & 4,18 & logika \\
         & DSP & 1 213 & 837 & 72 & 6 & 4,24 & logika \\
        \bottomrule
    \end{tabular}
\end{table}

% /TAB

\subsection{Architektura proti jazyku: histogram}

Histogram jsem ve VHDL navrhl ve dvou architekturách (tabulka~\ref{tab:histarch}). První ukládá 256 košů do bloku BRAM a~sousední pixely se stejnou hodnotou řeší předáváním mezivýsledku (forwarding). Druhá má 256 čítačů v~registrech. Obě dávají stejný výsledek, ale čítače zabírají 6 437 LUT místo 183, tedy 35krát víc. Volba architektury tu rozhoduje o~velikosti jádra víc než volba jazyka: verze v~HLS se stejnou architekturou zabírá 300 LUT.

% TAB tab:histarch

\begin{table}[htbp]
    \centering
    \caption{Histogram ve dvou architekturách (syntéza mimo kontext, cíl 5\,ns, jádro včetně AXI4-Lite)}
    \label{tab:histarch}
    \begin{tabular}{p{1.3cm}p{4.4cm}p{1.3cm}p{1.3cm}p{1.4cm}p{1.6cm}p{2.8cm}}
        \toprule
        Impl. & Architektura & LUT & FF & BRAM36 & Perioda [ns] & Omezení \\
        \midrule
        HLS & BRAM, předávání mezivýsledku & 300 & 422 & 1 & 3,08 & šířka pulzu BRAM \\
        VHDL & BRAM, předávání mezivýsledku & 183 & 288 & 1 & 3,08 & šířka pulzu BRAM \\
        VHDL & 256 čítačů v~registrech & 6 437 & 6 370 & 0,5 & 3,08 & šířka pulzu BRAM \\
        \bottomrule
    \end{tabular}
\end{table}

% /TAB

Perioda všech tří variant je 3,08\,ns a~omezuje ji minimální šířka pulzu hodin bloku BRAM. Varianta s~čítači blok BRAM v~jádře nemá, půl bloku zabírá okno pro čtení histogramu přes AXI4-Lite.

% TODO(ověřit): polovina BRAM36 u varianty s čítači – ověřit v utilization.rpt, že jde o okno hist v axil_regbank.

\subsection{Co hlásí nástroj a~co obvod zvládne}

Vitis HLS sám uvádí prostředky a~periodu po implementaci. Tabulka~\ref{tab:exportooc} je porovnává s~výsledkem stejného RTL ze syntézy mimo kontext. Počet LUT se liší nejvýš o~3\,\% (konvoluce 7×7 1 306 proti 1 341), počet FF nejvýš o~3.

% Pozn.: převod barev, převody RGB↔HSV a medián (sort3 změněno 19. 11.) v tabulce nejsou – export HLS jsem po vypnutí auto-rewind nespouštěl, starý export je z jiného RTL.  Perioda se liší až o 0,35 ns. U histogramu export uvádí 2,73 ns, ale blok BRAM vyžaduje periodu aspoň 3,077 ns, takže skutečné maximum je 325 MHz, ne 366 MHz. Export HLS totiž počítá jen s časovou rezervou logických cest a minimální šířku pulzu hodin nezahrnuje. Proto v celé kapitole uvádím jen čísla ze syntézy mimo kontext.

% TAB tab:exportooc

\begin{table}[htbp]
    \centering
    \caption{Stejné RTL z~HLS: výsledek exportu Vitis HLS proti syntéze mimo kontext stejným skriptem jako VHDL}
    \label{tab:exportooc}
    \begin{tabular}{p{3.0cm}p{2.4cm}p{2.4cm}p{1.9cm}p{1.9cm}p{2.8cm}}
        \toprule
        Jádro & LUT export / OOC & FF export / OOC & Perioda export [ns] & Perioda OOC [ns] & Omezení OOC \\
        \midrule
        histogram & 302 / 300 & 422 / 422 & 2,73 & 3,08 & šířka pulzu BRAM \\
        konvoluce 3×3 & 655 / 640 & 621 / 621 & 3,92 & 4,05 & logika \\
        konvoluce 5×5 & 959 / 950 & 863 / 861 & 4,17 & 4,49 & logika \\
        konvoluce 7×7 & 1 306 / 1 341 & 1 244 / 1 247 & 4,37 & 4,35 & logika \\
        \bottomrule
    \end{tabular}
\end{table}

% /TAB

Ještě větší rozdíl je mezi odhadem po syntéze HLS a~skutečností po implementaci. U~převodu HSV na RGB odhad uváděl 893 LUT, po implementaci jádro zabralo 556 LUT. Odhad syntézy HLS proto pro srovnání nepoužívám.

Rozptyl výsledků samotného rozmístění a~propojení jsem viděl po úpravě registrové banky, při které se obvod nezměnil (proměnnou jsem nahradil konstantou). U~25 z~28 konfigurací VHDL zůstala čísla stejná, u~tří jader s~registrovou bankou se perioda změnila o~0,2 až 1,0\,ns (konvoluce 3×3 bez DSP z~4,37 na 3,40\,ns). Rozdíly period pod přibližně 0,3\,ns proto nepovažuji za rozdíl mezi implementacemi.

\section{Propustnost na desce}
\label{sec:vdeska}

Propustnost jsem měřil na desce KV260 v~testovacím řetězci se všemi 14 jádry při 200\,MHz (oddíl~\ref{sec:deska}). Snímek 1920\,×\,1080 bodů v~RGB posílá VDMA z~paměti DDR přes jádro zpět do paměti. Doba snímku je průměr z~20 snímků, stav přenosu měřicí aplikace čte po 100\,µs, rozlišení měření je tedy 5\,µs na snímek. Výsledky z~měření s~aplikací přeloženou na -O2 jsou v~tabulce~\ref{tab:propustnost}, na úrovních -O0 a~-O3 se doba snímku lišila nejvýš o~1\,µs.

% TAB tab:propustnost

\begin{table}[htbp]
    \centering
    \caption{Doba zpracování snímku 1920\,×\,1080 na desce KV260 (200\,MHz, 20 snímků na měření, teoretický strop 10 368\,µs)}
    \label{tab:propustnost}
    \begin{tabular}{p{3.6cm}p{1.8cm}p{1.8cm}p{1.5cm}p{1.5cm}p{3.0cm}}
        \toprule
        Algoritmus & HLS [µs] & VHDL [µs] & HLS [snímků/s] & VHDL [snímků/s] & Proti průchodu HLS / VHDL [µs] \\
        \midrule
        průchod RGB & – & 10 371 & – & 96,4 & 0 \\
        průchod šedá & – & 10 371 & – & 96,4 & 0 \\
        převod barev RGB→YCbCr & 10 401 & 10 371 & 96,1 & 96,4 & +30 / 0 \\
        RGB→HSV & 10 412 & 10 371 & 96,0 & 96,4 & +41 / 0 \\
        HSV→RGB & 10 407 & 10 371 & 96,0 & 96,4 & +36 / 0 \\
        prahování (Otsu) & 10 371 & 10 371 & 96,4 & 96,4 & 0 / 0 \\
        medián 3×3 & 10 386 & 10 386 & 96,2 & 96,2 & +15 / +15 \\
        Sobel X 3×3 & 10 386 & 10 386 & 96,2 & 96,2 & +15 / +15 \\
        histogram & 10 371 & 10 371 & 96,4 & 96,4 & 0 / 0 \\
        \bottomrule
    \end{tabular}
\end{table}

% /TAB

Při zpracování 1 pixelu za takt trvá snímek 1920\,×\,1080 / 200\,MHz = 10 368\,µs. Průchod řetězcem bez jádra trval 10 371\,µs, tedy o~0,03\,\% déle. Všechna jádra pracují na tomto stropu. Okénková jádra (medián a~Sobel) přidávají 15\,µs, což je latence přibližně dvou řádků, než se naplní okno 3\,×\,3. Jádra HLS pro převod barev a~převod mezi RGB a~HSV jsou pomalejší o~30 až 41\,µs. To je cena za vypnutí auto-rewind (oddíl~\ref{sec:problemy}): mezi řádky přibude přibližně 5 taktů, u~převodu barev 5\,×\,1080 taktů = 27\,µs na snímek.

% TODO(odkaz): oddíl Nalezené problémy v kap. 5 má mít návěští sec:problemy (kap. 4 na něj odkazuje jako na 5.14).

Propustnost tedy jádra HLS a~VHDL nerozlišuje. Obě verze zpracují 1 pixel za takt a~rozdíl mezi nimi je nejvýš 0,4\,\% doby snímku. Rozdíl mezi verzemi je v~prostředcích a~maximální frekvenci (oddíl~\ref{sec:vooc}), ne v~propustnosti při 200\,MHz.

\section{Srovnání hardwaru a~softwaru}
\label{sec:vhwsw}

Softwarovou verzí je referenční model \texttt{lib\_imgproc} přeložený pro procesor Cortex-A53 na 1333\,MHz bez operačního systému, se zapnutou cache (oddíl~\ref{sec:sw}). Doba je medián z~10 běhů na snímku 1920\,×\,1080. Tabulka~\ref{tab:hwsw} uvádí dobu softwarového a~hardwarového zpracování a~jejich poměr. Hardwarový čas je u~jader VHDL, jádra HLS jsou pomalejší nejvýš o~0,4\,\%, což poměr na uvedenou přesnost nemění.

% TAB tab:hwsw

\begin{table}[htbp]
    \centering
    \caption{Doba zpracování snímku 1920\,×\,1080 v~softwaru na Cortex-A53 (1333\,MHz, bez operačního systému) a~v~hardwaru (VHDL, 200\,MHz) a~poměr SW/HW}
    \label{tab:hwsw}
    \begin{tabular}{p{3.4cm}p{1.4cm}p{1.4cm}p{1.4cm}p{1.3cm}p{1.3cm}p{1.3cm}p{1.3cm}}
        \toprule
        Algoritmus & SW -O0 [ms] & SW -O2 [ms] & SW -O3 [ms] & HW [ms] & SW/HW -O0 & -O2 & -O3 \\
        \midrule
        převod barev RGB→YCbCr & 1 147,7 & 278,0 & 228,7 & 10,37 & 111× & 27× & 22× \\
        RGB→HSV & 544,5 & 182,1 & 182,1 & 10,37 & 52× & 18× & 18× \\
        HSV→RGB & 521,6 & 115,8 & 115,8 & 10,37 & 50× & 11× & 11× \\
        prahování (Otsu) & 206,9 & 57,9 & 43,4 & 10,37 & 20× & 5,6× & 4,2× \\
        medián 3×3 & 2 728,1 & 549,9 & 561,1 & 10,39 & 263× & 53× & 54× \\
        Sobel X 3×3 & 2 334,1 & 567,4 & 538,5 & 10,39 & 225× & 55× & 52× \\
        histogram & 91,1 & 35,2 & 35,1 & 10,37 & 8,8× & 3,4× & 3,4× \\
        \bottomrule
    \end{tabular}
\end{table}

% /TAB

Při překladu s~-O2 je hardware rychlejší 3,4krát (histogram) až 55krát (Sobel). Velikost zrychlení odpovídá tomu, kolik práce připadá v~softwaru na jeden pixel. Histogram v~softwaru jen přičte jedničku do koše, takže procesor zpracuje snímek za 35,2\,ms. Medián a~Sobel počítají pro každý pixel okno 3\,×\,3, medián navíc třídí 9 hodnot, a~procesor na snímek potřebuje 550 a~567\,ms. Hardware zpracuje všechny algoritmy za stejnou dobu, protože každé jádro zpracuje 1 pixel za takt bez ohledu na to, kolik operací pro něj provede.

Vliv optimalizace překladače ukazuje tabulka~\ref{tab:prekladac}. Přechod z~-O0 na -O2 zrychlil software 2,6krát až 5,0krát. Přechod z~-O2 na -O3 pomohl jen převodu barev (o~17,7\,\%), prahování (o~25,0\,\%) a~Sobelovi (o~5,1\,\%), medián se zpomalil o~2,0\,\%. Pro srovnání s~hardwarem proto beru jako výchozí -O2 a~-O3 uvádím jako nejrychlejší software. Rozdíly pod 1\,\% mezi úrovněmi nepovažuji za skutečné.

% TAB tab:prekladac

\begin{table}[htbp]
    \centering
    \caption{Vliv optimalizace překladače na dobu softwarového zpracování}
    \label{tab:prekladac}
    \begin{tabular}{p{4.6cm}p{4.4cm}p{4.4cm}}
        \toprule
        Algoritmus & Zrychlení -O0 → -O2 & Změna -O2 → -O3 \\
        \midrule
        převod barev RGB→YCbCr & 4,1× & −17,7\,\% \\
        RGB→HSV & 3,0× & 0,0\,\% \\
        HSV→RGB & 4,5× & 0,0\,\% \\
        prahování (Otsu) & 3,6× & −25,0\,\% \\
        medián 3×3 & 5,0× & +2,0\,\% \\
        Sobel X 3×3 & 4,1× & −5,1\,\% \\
        histogram & 2,6× & −0,1\,\% \\
        \bottomrule
    \end{tabular}
\end{table}

% /TAB

% TODO(Karel): SW čas je z jednoho běhu aplikace na úroveň (medián z 10 běhů funkce). Chceš ještě 2 další běhy board_test pro rozptyl? Zabere to asi 15 minut s deskou.

\section{Využití obvodu celým řetězcem}
\label{sec:vretezec}

Syntéza mimo kontext měří jádro samostatně. V~testovacím řetězci je 14 jader spolu s~VDMA, přepínači AXI4-Stream a~propojením k~procesoru v~jednom obvodu. Tabulka~\ref{tab:retezec} uvádí využití obvodu a~časovou rezervu finálního bitstreamu, se kterým jsem měřil na desce (oddíl~\ref{sec:vdeska}).

% TAB tab:retezec

\begin{table}[htbp]
    \centering
    \caption{Využití obvodu xck26 testovacím řetězcem se 14 jádry a~časová rezerva při 200\,MHz (finální bitstream, Vivado 2025.2)}
    \label{tab:retezec}
    \begin{tabular}{p{2.2cm}p{1.6cm}p{1.6cm}p{2.3cm}p{2.3cm}p{2.0cm}p{1.9cm}p{1.6cm}}
        \toprule
         & WNS [ns] & WHS [ns] & LUT & FF & BRAM36 & DSP & URAM \\
        \midrule
        použito & +0,229 & +0,010 & 14 006 (12,0\,\%) & 17 969 (7,7\,\%) & 16 (11,1\,\%) & 64 (5,1\,\%) & 2 (3,1\,\%) \\
        k~dispozici &  &  & 117 120 & 234 240 & 144 & 1 248 & 64 \\
        \bottomrule
    \end{tabular}
\end{table}

% /TAB

Řetězec zabírá 12\,\% buněk LUT a~11\,\% bloků BRAM obvodu. Jak se prostředky dělí mezi jednotlivé bloky, ukazuje tabulka~\ref{tab:retezecbloky}. Všech 14 jader zabírá dohromady 37\,\% LUT řetězce, zbytek připadá na přenos dat: VDMA, přepínače a~propojení k~procesoru a~do paměti DDR. Omezením tedy není velikost obvodu, ale časová rezerva.

% TAB tab:retezecbloky

\begin{table}[htbp]
    \centering
    \caption{Rozdělení prostředků testovacího řetězce po blocích (finální bitstream, URAM jen ve VDMA)}
    \label{tab:retezecbloky}
    \begin{tabular}{p{4.6cm}p{1.6cm}p{1.6cm}p{1.6cm}p{1.6cm}p{1.6cm}}
        \toprule
        Blok & LUT & Podíl LUT & FF & BRAM36 & DSP \\
        \midrule
        7 jader HLS & 2 647 & 19\,\% & 2 450 & 6 & 32 \\
        7 jader VHDL & 2 525 & 18\,\% & 2 585 & 5 & 26 \\
        převod na šedou a~zpět (VHDL) & 371 & 3\,\% & 557 & 0 & 6 \\
        VDMA & 2 973 & 21\,\% & 4 636 & 5 & 0 \\
        přepínače AXI4-Stream & 1 436 & 10\,\% & 1 899 & 0 & 0 \\
        propojení AXI4-Lite (řízení) & 1 962 & 14\,\% & 1 882 & 0 & 0 \\
        propojení do paměti DDR & 2 097 & 15\,\% & 3 911 & 0 & 0 \\
        \bottomrule
    \end{tabular}
\end{table}

% /TAB

Časová rezerva celého řetězce je při 200\,MHz 0,229\,ns. Kritická cesta leží uvnitř jádra HLS pro převod HSV na RGB, v~násobení saturace a~zlomkové části odstínu. Vivado ho umístil do bloku DSP včetně předsčítačky, ale bez vnitřních registrů. Cesta trvá 4,63\,ns a~prochází 11 úrovněmi logiky, z~toho 8 uvnitř bloku DSP. Při syntéze mimo kontext měly obě varianty jádra (s~násobením v~LUT i~v~DSP) rezervu 0,77 až 0,95\,ns, v~řetězci přibylo směrování a~rozmístění vedle ostatních jader.

Tohle je příklad, jak se chová jádro HLS navržené na hranu zadané periody. Při 5\,ns vloží HLS mezi registry tolik logiky, kolik se do periody vejde, a~v~hustším návrhu pak rezerva mizí. Jádra ve VHDL s~registrem za každou operací kritickou cestu řetězce netvořila. Nápravou by byla varianta s~latencí násobení 3 takty z~tabulky~\ref{tab:hsvdsp}.

% TODO(Karel): V řetězci zabírá jádro hls_hsv2rgb 232 LUT, 264 FF a 11 DSP, tedy jako varianta s násobením v DSP (OOC 224 LUT, 241 FF, 11 DSP), ne jako výchozí varianta (OOC 556 LUT, 1 DSP). Možné příčiny: (a) všechny tři varianty HSV→RGB mají stejný název IP xilinx.com:hls:hls_hsv2rgb:1.0 a build_harness.tcl přidává do katalogu všechny tři adresáře, Vivado tedy mohl vzít variantu _dsp; (b) syntéza IP v blokovém návrhu sama mapuje násobení 8 × 8 do DSP (sekce 12 z 22. 10., kdy varianta _dsp ještě neexistovala, už ukazovala kritickou cestu v DSP). Do vyjasnění v textu netvrdím, která varianta v řetězci je. Oprava: dát variantám různé názvy IP (package.ip.name) a bitstream přeložit znovu.

\section{Diskuse}
\label{sec:diskuse}

% TODO(text): diskusi napíšu po kontrole čísel a tvém schválení tónu. Body podle osnovy a podkladů:

% 1. Propustnost HLS a VHDL nerozlišuje: obě 1 px/takt, rozdíl ≤ 0,4 % (cena vypnutí auto-rewind).

% 2. VHDL: kratší nebo stejná perioda u všech 9 jader, menší LUT u 7 z 9; HLS: méně FF u převodů RGB↔HSV, rychlejší vývoj (doplnit tvou zkušenost – kolik času zabralo jádro v HLS a ve VHDL, pokud to víš).

% 3. Architektura rozhoduje víc než jazyk (histogram 35× LUT).

% 4. Jeden parametr LUT ↔ DSP: ve VHDL stačí jeden, v HLS dva a je potřeba vědět o latenci.

% 5. Číslo z nástroje nestačí: export HLS bez šířky pulzu, odhad syntézy HLS +61 % LUT, rozptyl P&R až 1 ns.

% 6. HW proti SW: 3,4× až 55× při -O2; zrychlení roste s počtem operací na pixel, HW čas je pro všechny stejný.

% 7. Limity měření: jeden běh SW na úroveň, příkon neměřen, na desce jen Sobel X, zpětný převod YCbCr jen v simulaci, finální využití řetězce doplnit.
